\section{Gamma identities, harmonic values, and retained primitives}
\label{sec:gamma}

The normal-form theorem is an identity in a free differential algebra.
Its analytic meaning must keep both the primitive in
\eqref{eq:decomp} and the distinction between a projection and the full
normal polynomial.

\begin{theorem}[Transport with the exact endpoint]\label{thm:endpoint}
Let $0<A<B$, $m\ge0$, and write
$U(x)=(\psi(x),\psi'(x),\ldots)$. Then
\begin{equation}\label{eq:Gamma-endpoint}
\begin{split}
 \int_A^B N_m(U(x))\dd x
 ={}&(-1)^m\int_A^B\Gamma(x)
       \frac{\mathrm d^m}{\mathrm dx^m}\frac1{\Gamma(x)}\dd x\\
 &-H_m(U(B))+H_m(U(A)).
\end{split}
\end{equation}
The polynomial $H_m$ is computed by the terminating rational construction
in Proposition~\ref{prop:normal}. In particular, the formula gives an
anchored primitive for every normal combination $N_m$ without suppressing
an integration-by-parts boundary constant.
\end{theorem}
\begin{proof}
Substitute $u_r=\psi^{(r)}(x)$ into
$A_m=N_m+\D H_m$. Under this substitution $\D$ is ordinary argument
differentiation. Use \eqref{eq:Gamma-realization} and integrate over the
compact interval. All functions are analytic there.
\end{proof}

For use with the compact cubic formula, let
$R_m=N_m-\sigma_3N_m$. The precise projected identity is therefore
\begin{equation}\label{eq:projected-endpoint}
\begin{split}
 \int_A^B(\sigma_3N_m)(U(x))\dd x
 ={}&(-1)^m\int_A^B\Gamma(x)(\Gamma^{-1})^{(m)}(x)\dd x\\
 &-[H_m(U(x))]_A^B-\int_A^BR_m(U(x))\dd x.
\end{split}
\end{equation}
One must not erase the last term when the higher-derivative sector is
nonzero. For $m\le9$ it vanishes by weight, but it need not vanish from
weight ten onward.

\subsection{Two explicit normal polynomials}
The weight-eight identity is
\begin{equation}\label{eq:N8}
 N_8=u_0^8-70u_0^4u_1^2+28u_0^2u_2^2-7u_1^4-u_3^2.
\end{equation}
For this case the complete endpoint polynomial is
\begin{equation}\label{eq:H8}
\begin{split}
 H_8={}&-4u_0^7+56u_0^5u_1-70u_0^4u_2-140u_0^3u_1^2
       +56u_0^3u_3\\
 &+252u_0^2u_1u_2-28u_0^2u_4+112u_0u_1^3
       -112u_0u_1u_3\\
 &-84u_0u_2^2+8u_0u_5-98u_1^2u_2+20u_1u_4+36u_2u_3-u_6.
\end{split}
\end{equation}
Direct differentiation verifies $A_8=N_8+\D H_8$.
Equations~\eqref{eq:Gamma-endpoint}, \eqref{eq:N8}, and \eqref{eq:H8}
are an ordinary finite-interval integral identity involving digamma,
trigamma, and higher polygammas, with an entirely explicit endpoint.

At weight eleven,
\begin{equation}\label{eq:N11}
\begin{split}
 N_{11}={}&u_0^{11}-330u_0^7u_1^2+462u_0^5u_2^2
 -1155u_0^3u_1^4-165u_0^3u_3^2\\
 &+550u_0^2u_2^3+990u_0u_1^2u_2^2+11u_0u_4^2-55u_2u_3^2.
\end{split}
\end{equation}
Thus $R_{11}=11u_0u_4^2$. This gives a concrete test of why the
projection alone cannot be substituted for the full $N_{11}$ in an
unqualified Gamma integral. The coefficient $-55$ is the first coupled
case of \eqref{eq:square-answer}. The full $H_{11}$ and other exact
primitives are supplied as rational JSON certificates.

\subsection{Generalized harmonic numbers}
Let $H_N^{(r)}=\sum_{j=1}^Nj^{-r}$, $H_N=H_N^{(1)}$, and let
$\gamma$ denote Euler's constant. The standard shift formulas give
\begin{equation}\label{eq:harmonic-substitution}
 u_0=H_N-\gamma,\qquad
 u_r=(-1)^{r+1}r!\bigl(\zeta(r+1)-H_N^{(r+1)}\bigr),\quad r\ge1,
\end{equation}
at $x=N+1$ \cite{DLMF5,DLMF25}. Let $U_N$ denote this tuple. Then
\begin{equation}\label{eq:harmonic-identity}
 N_m(U_N)+(\D H_m)(U_N)
 =m![t^m]\frac{\Gamma(N+1)}{\Gamma(N+1-t)}.
\end{equation}
The right side also equals
\begin{equation}\label{eq:harmonic-generator}
 m![t^m]\exp\left((H_N-\gamma)t-
 \sum_{r\ge2}\frac{\zeta(r)-H_N^{(r)}}r\,t^r\right).
\end{equation}
This follows by substituting \eqref{eq:harmonic-substitution} into
\eqref{eq:Bell}. It is initially a local analytic series in $t$ and an
identity of every finite coefficient. The derivative $\D H_m$ is formed
in the free differential algebra \emph{before} substitution; it is not a
finite difference in $N$.

Consequently \eqref{eq:square-answer} and the nonlinear coefficients
specify exact mixed harmonic-tail contributions to every identity
\eqref{eq:harmonic-identity}. This is a systematic family of identities,
not a claim that the individual harmonic products have been shown
arithmetically independent or that every integral of one isolated
polygamma monomial has been separately evaluated.
